\section{人工智能工具使用说明}\label{app:ai-usage}

\subsection{工具信息与使用范围}
本队在解题过程中使用OpenAI Codex作为辅助工具。使用信息如下：工具名称为OpenAI Codex桌面端，开发机构为OpenAI；所用模型属于GPT-5系列Codex模型。客户端未向参赛队显示每次会话所对应的精确底层快照，因此本文不虚构子版本号。公开模型资料将GPT-5-Codex说明为针对Codex智能编程任务优化、且底层快照会随平台更新的GPT-5版本；本队实际使用日期为2026年9月24日至26日，官方资料查询地址为\url{https://developers.openai.com/api/docs/models/gpt-5-codex}。

人工智能工具参与了信息检索、赛题理解与任务拆解、建模思路讨论、LaTeX排版、辅助编程、数据质量检查、图表制作和文字初稿整理。它不直接决定最终模型与结论，所有输出均作为候选材料，由队员结合题目附件、程序实算结果和约束检查决定是否采用。

\subsection{代表性输入内容与后续处理}
表\ref{tab:ai-usage}汇总主要使用场景。表中“输入内容”是对多轮交互任务的归纳，不包含队伍编号、个人信息或其他与解题无关的内容。

\begin{table}[H]
\centering
\caption{人工智能工具的主要使用场景与处理策略}\label{tab:ai-usage}
\small
\begin{tabularx}{\textwidth}{>{\raggedright\arraybackslash}p{2.6cm}YY}
\toprule
使用环节 & 代表性输入内容 & 队员的核验与后续处理\\
\midrule
信息检索 & 检索无人机载荷相关路径、带时间窗邻域搜索、空地通信和图划分的正式文献 & 仅保留可由出版社、DOI或标准组织页面追溯的文献；核对作者、题名、年份和适用范围后再引用\\
赛题理解与任务拆解 & 梳理四问的输入、输出、约束、依赖关系及容易混淆的时间口径 & 队员回到题目和附件逐项核对，并确定四问共用地形、能量、资源事件时间轴和结果字段\\
建模思路讨论 & 比较单点组批、异构调度、连续通信认证、中继排程和固定方案分区的候选方法 & 队员讨论模型假设、评价指标与求解边界；对未穷尽的搜索明确只报告可行方案，不作全局最优声明\\
辅助编程与调试 & 生成或修改数据读取、DEM航段遍历、调度、通信验证、分区枚举和结果导出代码 & 在本地运行程序，以80箱唯一覆盖、硬时限、返航SOC、资源无冲突和连续通信等约束逐项检查；结果冲突时以题目原始数据和重新计算为准\\
数据质量检查 & 对比CSV、JSON、图表和正文中的架次数、能耗、完工时间、通信分段及资源缺口 & 通过汇总回算、跨文件交叉检查和多轮复核定位差异；未经复核的数字不写入结论\\
图表与论文写作 & 生成流程图、结果图、LaTeX段落和参考文献格式的候选稿 & 队员逐问审读并提出修改，检查图文对应、术语和数值一致性，删除无法确认的表述后形成提交稿\\
\bottomrule
\end{tabularx}
\end{table}

\subsection{队员的独立判断与责任}
队员在使用人工智能工具前后持续参与问题判断和方案取舍。具体包括：根据题意确定完整货箱不可拆、返航安全余量和硬时限优先；比较问题二基线与22架次方案后选择后者；针对问题三与其他队伍结果的差异，主动质疑当前方案的最优性，并决定在本文中保留已经完整验证的可行方案，同时明确其搜索边界而不夸大全局最优性；在问题四中确认固定问题三任务及时序后再评价独立分组的资源缺口；对生成的PDF逐页检查并要求补充模型检验、误差分析、结论、文献角标和四问流程图。

最终模型选择、参数口径、结果采用和论文提交均由参赛队负责。人工智能输出不能替代题目附件、正式文献和可复现计算；对于工具给出的建议，只有在队员能够解释其含义、确认数据来源并通过程序或逻辑复核后才予以采用。
